% Einheit 1 von 2 – Kollinearität und lineare Abhängigkeit (bank/linearkombination-und-lineare-abhaengigkeit/e1.jsonl, zusammenbau v0.1)
%% TODO Zeitmarke Einheit 1: Marken-Zeile nicht lesbar
%% TODO Zweigzeile Teil 1: was hier gelernt wird (Satz in Schülersprache aus dem Einheitstitel) – Einheit 1
\einheitenkopf[e1]{Einheit 1 von 2 · Kollinearität und lineare Abhängigkeit}
\zweigzeile{keine Prüfungsaufgabe}
\setcounter{aufgabe}{6}

%% TODO Titel: Ich-kann-Satz fehlt in der Bank (Platzhalter: Kettenname) – Nr. 7
%% TODO Anweisung: ein Satz über der ersten Teilaufgabe fehlt in der Bank – Nr. 7
\begin{aufgabe}{Kollinearität und lineare Abhängigkeit}
%% TODO \rechenplatz{n}: Schrittzahl des Grundfalls fehlt in der Bank (nur bei fester Schreibform, 2.3 b)
\begin{teile}
\teil $\vec{a} = (1 | -2 | 3)$ und $\vec{b} = (-2 | 4 | -6)$. Ist einer ein Vielfaches des anderen? \\ \kreuz{ja} \kreuz{nein}
\teil Sind $\vec{u} = (1 | 3 | -1)$ und $\vec{v} = (4 | 12 | -4)$ kollinear? Gib gegebenenfalls den Faktor an.
\teil Sind $\vec{u} = (2 | -3 | 1)$ und $\vec{v} = (-4 | 6 | -2)$ kollinear? Zeigen sie in dieselbe Richtung oder in die Gegenrichtung?
\teil Sind $\vec{a} = (1 | 0 | 1)$, $\vec{b} = (0 | 1 | 1)$ und $\vec{c} = (2 | 3 | 5)$ linear abhängig?
\teil Weise nach, dass das Viereck $ABCD$ mit $A(1 | 1 | 0)$, $B(7 | 4 | 3)$, $C(4 | 4 | 2)$ und $D(2 | 3 | 1)$ ein Trapez ist.
\end{teile}
\end{aufgabe}
